\documentclass[14pt,a4paper,oneside]{extbook}
\usepackage[utf8]{vietnam}
\usepackage{amsmath,amssymb,mathrsfs,fancyhdr,enumerate,multirow,makecell}
\usepackage{tkz-euclide,tikz-3dplot,tikz,tkz-tab}
\usetikzlibrary{arrows.meta,calc,intersections,angles,quotes,backgrounds,shapes.geometric,patterns,through,shadows,positioning}
\usepackage{yhmath}
\usepackage[top=0.7cm, bottom=0.7cm, left=1.0cm, right=1.0cm, headheight=16pt, headsep=3mm, footskip=5mm]{geometry}
\usepackage[hidelinks,unicode]{hyperref}
\usepackage{qrcode}
\usepackage[many]{tcolorbox}
\usepackage{fontawesome}

\pagestyle{fancy}
\fancyhf{}
\fancyhead[L]{\footnotesize\bfseries\color{teal!85!black}\faBook~SỔ TAY KIẾN THỨC TOÁN 4}
\fancyhead[R]{\footnotesize\bfseries\color{blue!85!black}Lớp Toán Thầy Lĩnh -- loptoanthaylinh.vn}
\fancyfoot[C]{\footnotesize\bfseries\color{gray!80!black}-- Sổ tay kiến thức: Bài 29. Hai đường thẳng song song --}
\renewcommand{\headrulewidth}{0.4pt}

\begin{document}

% --- BANNER TIÊU ĐỀ BÀI HỌC ---
\noindent
\begin{tcolorbox}[
  enhanced,
  colback=teal!8!white,
  colframe=teal!80!black,
  arc=3mm, boxrule=1pt,
  left=3mm, right=3mm, top=2mm, bottom=2mm,
  halign=center
]
  {\Large\bfseries\color{teal!90!black}\faBookmark\ BÀI 29: HAI ĐƯỜNG THẲNG SONG SONG}\\[1.5pt]
  {\footnotesize\itshape\color{black!75} Khái niệm không cắt nhau, khoảng cách không đổi và dấu hiệu nhận biết qua đường trung gian trên 1 trang A4}
\end{tcolorbox}

\vspace{2mm}

% --- KHỐI 1 & 2: LÝ THUYẾT & KĨ NĂNG ---
\noindent
\begin{minipage}[t]{0.49\textwidth}
  \begin{tcolorbox}[
    enhanced,
    title={\footnotesize\bfseries\color{white}\faArrowsH\ 1. Khái niệm hai đường thẳng song song},
    colback=blue!2!white,
    colframe=blue!75!black,
    arc=2.5mm, boxrule=0.8pt, top=1.8mm, bottom=1.8mm, left=2mm, right=2mm
  ]
    \footnotesize
    $\bullet$ \textbf{Hai đường thẳng song song với nhau thì không bao giờ cắt nhau} (dù kéo dài mãi mãi về hai phía).\\
    $\bullet$ Kí hiệu: $a \parallel b$ (đường thẳng $a$ song song với đường thẳng $b$).\\
    $\bullet$ \textbf{Hình ảnh thực tế}:\\
    --- Hai thanh ray đường tàu hỏa.\\
    --- Hai mép đối diện của bảng, cánh cửa sổ.\\
    --- Các dòng kẻ ngang trên trang vở ô li.
  \end{tcolorbox}
\end{minipage}%
\hfill
\begin{minipage}[t]{0.49\textwidth}
  \begin{tcolorbox}[
    enhanced,
    title={\footnotesize\bfseries\color{white}\faExchange\ 2. Dấu hiệu nhận biết quan trọng},
    colback=teal!2!white,
    colframe=teal!75!black,
    arc=2.5mm, boxrule=0.8pt, top=1.8mm, bottom=1.8mm, left=2mm, right=2mm
  ]
    \footnotesize
    $\bullet$ \textbf{Quan hệ với đường thẳng vuông góc}:\\
    Nếu hai đường thẳng cùng \textbf{vuông góc} với một đường thẳng thứ ba thì chúng \textbf{song song} với nhau.
    \begin{center}
      \textbf{Nếu $a \perp c$ và $b \perp c$ thì $a \parallel b$}
    \end{center}
    $\bullet$ \textbf{Khoảng cách giữa hai đường song song}: Luôn luôn bằng nhau tại mọi vị trí đo.
  \end{tcolorbox}
\end{minipage}

\vspace{2mm}

% --- KHỐI 3: SƠ ĐỒ TƯ DUY TIKZ ---
\noindent
\begin{tcolorbox}[
  enhanced,
  title={\footnotesize\bfseries\color{white}\faSitemap\ SƠ ĐỒ TƯ DUY: BẢN CHẤT HÌNH HỌC CỦA HAI ĐƯỜNG THẲNG SONG SONG},
  colback=yellow!2!white,
  colframe=orange!85!black,
  arc=2.5mm, boxrule=0.8pt, top=1.5mm, bottom=1.5mm, left=1mm, right=1mm,
  halign=center
]
  \begin{tikzpicture}[
    scale=0.92,
    box/.style={rectangle, rounded corners=2mm, draw=teal!80!black, fill=teal!10, thick, inner sep=3pt, font=\scriptsize\bfseries, align=center},
    arr/.style={-{Stealth[scale=0.85]}, thick, draw=black!70}
  ]
    % 2 đường thẳng song song
    \draw[very thick, blue!85!black] (0.5, 1.0) -- (6.5, 1.0) node[right, font=\small\bfseries] {$a$};
    \draw[very thick, blue!85!black] (0.5, -0.6) -- (6.5, -0.6) node[right, font=\small\bfseries] {$b$};

    % Đoạn vuông góc trung gian thể hiện khoảng cách
    \draw[thick, dashed, red!80!black] (2.0, -0.6) -- (2.0, 1.0) node[midway, left, font=\tiny\bfseries\color{red!80!black}]{$d$};
    \draw[thick, dashed, red!80!black] (5.0, -0.6) -- (5.0, 1.0) node[midway, left, font=\tiny\bfseries\color{red!80!black}]{$d$};

    % Kí hiệu vuông góc
    \draw[thick, red!80!black] (2.0, 0.8) -- (2.2, 0.8) -- (2.2, 1.0);
    \draw[thick, red!80!black] (2.0, -0.4) -- (2.2, -0.4) -- (2.2, -0.6);
    \draw[thick, red!80!black] (5.0, 0.8) -- (5.2, 0.8) -- (5.2, 1.0);
    \draw[thick, red!80!black] (5.0, -0.4) -- (5.2, -0.4) -- (5.2, -0.6);

    % Chú thích bên phải
    \node[box, draw=blue!75!black, fill=blue!8, text width=7.5cm, inner sep=2.5pt] (c1) at (11.5, 0.8) {\textbf{\color{blue!85!black}KHÔNG BAO GIỜ CẮT NHAU}\\Kéo dài vô tận về 2 phía đều không có điểm chung\\$a \parallel b$};
    \node[box, draw=teal!80!black, fill=teal!8, text width=7.5cm, inner sep=2.5pt] (c2) at (11.5, -0.8) {\textbf{\color{teal!85!black}KHOẢNG CÁCH KHÔNG ĐỔI}\\Cùng vuông góc với một đường thẳng\\$\rightarrow$ Khoảng cách $d$ bằng nhau ở mọi vị trí};

    \draw[arr] (5.5, 0.3) -- (c1.west);
    \draw[arr] (5.5, -0.3) -- (c2.west);
  \end{tikzpicture}
\end{tcolorbox}

\vspace{2mm}

% --- KHỐI 4: VÍ DỤ MINH HỌA & BẪY TRẮC NGHIỆM ---
\noindent
\begin{minipage}[t]{0.53\textwidth}
  \begin{tcolorbox}[
    enhanced,
    title={\footnotesize\bfseries\color{white}\faCheckCircle\ Ví dụ thực hành chuẩn mực},
    colback=green!2!white,
    colframe=green!60!black,
    arc=2.5mm, boxrule=0.8pt, top=1.8mm, bottom=1.8mm, left=2.2mm, right=2.2mm
  ]
    \footnotesize
    Trong hình chữ nhật $ABCD$:\\
    $\bullet$ Cặp cạnh $AB$ và $CD$ song song với nhau ($AB \parallel CD$).\\
    $\bullet$ Cặp cạnh $AD$ và $BC$ song song với nhau ($AD \parallel BC$).\\
    $\Rightarrow$ Hình chữ nhật có \textbf{2 cặp cạnh đối diện song song} và bằng nhau từng đôi một: $AB = CD$ và $AD = BC$.
  \end{tcolorbox}
\end{minipage}%
\hfill
\begin{minipage}[t]{0.45\textwidth}
  \begin{tcolorbox}[
    enhanced,
    title={\footnotesize\bfseries\color{white}\faExclamationTriangle\ Mẹo tránh bẫy trắc nghiệm},
    colback=red!2!white,
    colframe=red!70!black,
    arc=2.5mm, boxrule=0.8pt, top=1.8mm, bottom=1.8mm, left=2.2mm, right=2.2mm
  ]
    \footnotesize
    $\bullet$ \textbf{Bẫy chưa cắt nhau trên hình vẽ:}\\
    Hai đoạn thẳng vẽ trên trang giấy chưa chạm vào nhau nhưng nếu \textbf{kéo dài ra mà cắt nhau} thì \textbf{KHÔNG PHẢI} hai đường thẳng song song!\\
    $\bullet$ Luôn kiểm tra xem khoảng cách giữa chúng có bằng nhau không.
  \end{tcolorbox}
\end{minipage}

\end{document}
